\section{Rozdział 4: Implementacja}

\subsection{Kluczowe metody klasy DoublyLinkedList}
Poniżej przedstawiono implementację metody dodawania elementu pod wskazany indeks z zachowaniem poprawności wskaźników \texttt{next} oraz \texttt{prev}:

\begin{lstlisting}[caption={Dodawanie elementu pod podany indeks}]
void addAt(int index, const T& val) {
    if (index < 0 || index > count) {
        throw std::out_of_range("Indeks poza zakresem!");
    }
    if (index == 0) {
        addFront(val);
    } else if (index == count) {
        addBack(val);
    } else {
        Node<T>* current = getNodeAt(index);
        Node<T>* newNode = new Node<T>(val);
        Node<T>* prevNode = current->prev;

        prevNode->next = newNode;
        newNode->prev = prevNode;
        newNode->next = current;
        current->prev = newNode;
        count++;
    }
}
\end{lstlisting}

\subsection{Repozytorium Git oraz historia commitów}
\begin{itemize}
    \item \textbf{Adres repozytorium GitHub:} \url{https://github.com/KrzysztofZdzieszynski/lista-dwukierunkowa}
\end{itemize}

W ramach pracy z systemem kontroli wersji rejestrowano postępy w repozytorium. Poniżej przedstawiono pełny log z historii commitów (\texttt{git log --oneline --graph --all}):

\begin{lstlisting}[caption={Rzeczywista historia commitów z systemu Git}]
* da16ced (HEAD -> main, origin/main) Commit 5: Dodanie wzorca Factory, wyswietlania wstecz oraz pliku main.cpp
* 2b3b477 Commit 4: Implementacja wzorca Iteratora oraz integracja z klasa listy
* 43acabd Commit 3: Dodanie metod removeFront, removeBack oraz czyszczenia listy
* 181c786 Commit 2: Dodanie metod addFront, addBack oraz wyswietlania listy
* e3a06e8 Commit 1: Dodanie struktury Node i bazowego szkieletu klasy listy
\end{lstlisting}

\subsection{Weryfikacja scenariuszy w Git}
\begin{enumerate}
    \item Utworzono 5 chronologicznych commitów dokumentujących postęp prac.
    \item Przeprowadzono testowy powrót w historii o dwa commity wstecz za pomocą polecenia \texttt{git reset --hard HEAD}.
    \item Sklonowano repozytorium do drugiego katalogu na dysku, wprowadzono zmiany, wysłano je na serwer zdalny (\texttt{git push}) i zsynchronizowano w pierwotnej lokalizacji przy użyciu \texttt{git pull}.
    \item Celowo usunięto plik \texttt{ListFactory.h}, wywołując błąd kompilacji, a następnie przywrócono go poleceniem \texttt{git checkout ListFactory.h}.
\end{enumerate}

\subsection{Generowanie dokumentacji Doxygen}
Plik konfiguracyjny \texttt{Doxyfile} został wygenerowany poleceniem \texttt{doxygen -g Doxyfile} z włączoną opcją \texttt{GENERATE\_LATEX = YES}. Wynikowe pliki trafiły do katalogu \texttt{latex/}, po czym zostały skompilowane poleceniem \texttt{pdflatex refman.tex} do pliku \texttt{refman.pdf}.